% !PRINTED 412
\weberpage{412}
%<<CONT
in $A$ noch acht cyklische Permutationen von je drei Ziffern
$\gamma = (0, 1, 2) \ldots$ vor, und $A$ kann so dargestellt werden:
\[
A = B + B\gamma + B\gamma^2.
\]
Je eine Permutation $\alpha$ und eine Permutation $\gamma$ können als
Erzeugende der Gruppe betrachtet werden. Denn ist etwa
\[
\gamma = (0, 1, 2),\ \alpha_1 = (0, 1)(2, 3),
\]
so ist
\weq{1}{\alpha_1\gamma\alpha_1 = \gamma' = (0, 3, 1),}
und aus $(0, 1, 2), (0, 3, 1)$ kann die ganze Gruppe $A$ abgeleitet
werden (Bd. I, \S.\,153, 7.).

Ebenso kann man $\gamma$ und $\gamma'$ als erzeugende Permutationen
von $A$ auffassen. Insbesondere ist
\weq{2}{\gamma'\gamma\gamma' = \alpha_1.}
Wir wählen $\vartheta$ so, dass der Permutation $\gamma$ eine Multiplication
$\varepsilon\vartheta$ entspricht, in der, da $\gamma$ vom Grade 3 ist, $\varepsilon$ eine dritte Einheitswurzel
 bedeutet. Wir setzen
\weq{3}{\vartheta = \frac{\varphi(x_0, x_1, x_2, x_3)}{\psi(x_0, x_1, x_2, x_3)},}
und verstehen unter $\varphi, \psi$ zwei ganze Functionen der vier
Variablen $x$ ohne gemeinschaftlichen Theiler.

Aus der identischen Gleichung
\[
\frac{\varphi_\gamma}{\psi_\gamma} = \varepsilon\,\frac{\varphi}{\psi}
\]
folgt dann, dass
\weq{4}{\varphi_\gamma = \varepsilon_1\varphi,\ \psi_\gamma = \varepsilon_2\varphi} % [sic]
sein muss, worin $\varepsilon_1, \varepsilon_2$ gleichfalls dritte Einheitswurzeln sind.

Wir wenden eine der Permutationen $\alpha$ auf $\vartheta$ an, und erhalten
 aus $\vartheta_\alpha = \chi(\vartheta)$:
\weq{5}{\frac{\varphi_\alpha}{\psi_\alpha} = \frac{a\varphi + b\psi}{c\varphi + d\psi}.}
Da nun $\varphi, \psi$, also auch $a\varphi + b\psi$ und $c\varphi + d\psi$, und
ebenso $\varphi_\alpha, \psi_\alpha$ ohne gemeinsamen Theiler sind, so muss in der
Identität (5) der Zähler dem Zähler und der Nenner dem
Nenner, wenigstens bis auf einen constanten Factor, gleich sein.
Diesen constanten Factor können wir in die Constanten $a, b, c, d$
einrechnen und erhalten
\[
\text{6)}\quad \varphi_\alpha = a\varphi + b\psi,\ \psi_\alpha = c\varphi + d\psi.
\]
